\chapter{Writing and Typesetting Formal Logic}
\label{ch:latex-style}

\section{Academic prose carries the argument}
\label{sec:latex-style-01}

Every display belongs to a grammatical sentence.  Use prose to state why an
equation or derivation appears, which hypotheses are active, and what the
display establishes.  Symbols should compress repeated structure after it has
been defined; they should not replace the definition.

\begin{method}[Paragraph structure]
A proof paragraph should normally have four components:
\begin{enumerate}
  \item the local target;
  \item the definition or prior result used;
  \item the derivation, construction, or calculation; and
  \item the exact consequence for the theorem.
\end{enumerate}
When the calculation has several lines, align relation symbols and justify each
nontrivial transition.
\end{method}

The following rules govern academic style in this reference.
\begin{enumerate}
  \item State the language, calculus, semantic class, and ambient metatheory
        before asserting derivability, validity, completeness, or consistency.
  \item Quantify every object on which a claim depends.  Replace ``for suitable
        objects'' by the exact domain and hypotheses.
  \item Use \emph{definition}, \emph{calculation}, \emph{derivation}, and
        \emph{countermodel} only for the corresponding evidence.  A diagram or
        suggestive example is not a derivation.
  \item Use only the proof-status descriptions \emph{complete proof},
        \emph{proof sketch}, and \emph{cited theorem}.  A proof sketch must
        state its omitted obligations; a cited theorem must reproduce the
        hypotheses used and identify its source.
  \item Replace evaluative adverbs such as ``clearly'' and ``obviously'' by the
        definition, rule instance, or finite calculation that licenses the
        sentence.
  \item Keep object-language formulas distinct from metalinguistic assertions.
        In particular, \(\varphi\leftrightarrow\psi\) is a formula, whereas
        \(\varphi\equiv\psi\) is a declared metalinguistic relation.
  \item State the induction object, induction hypothesis, and every constructor
        case.  If a simultaneous or strengthened induction is required, state
        it before the base case.
  \item End a proof paragraph with the exact clause established.  Do not leave
        the reader to infer whether the argument proved existence, uniqueness,
        equivalence, or only one implication.
\end{enumerate}

